\section{Method}
\label{sec:method}

\Cref{fig:overview} gives the overview of the method, from its input to its output and the feedback that closes the loop.
Open the section with the modelling choice, then give each component its motivation, its design and the reason for that design.

\begin{figure*}[t]
  \centering
  % Overview of the method: input, method, output and a feedback loop.
\begin{tikzpicture}[
    font=\sffamily\small,
    stage/.style={rounded corners=5pt, minimum height=1.25cm, minimum width=3.3cm,
      align=center, draw=milapurple!55, line width=0.6pt, fill=labtint},
    core/.style={stage, fill=milapurple, draw=milapurple, text=white},
    flow/.style={-latex, line width=0.8pt, draw=milapurple!70}]
  \node[stage] (in)  at (0,0)   {\textbf{Input}\\[1pt]\footnotesize task and context};
  \node[core]  (m)   at (4.6,0) {\textbf{\method{}}\\[1pt]\footnotesize the component that is new};
  \node[stage] (out) at (9.2,0) {\textbf{Output}\\[1pt]\footnotesize answer and trace};
  \node[stage, draw=udemblue!60, fill=udemblue!7, minimum width=4.2cm]
    (fb) at (4.6,-1.85) {\textbf{Feedback}\\[1pt]\footnotesize reward, critique or memory};
  \draw[flow] (in) -- (m);
  \draw[flow] (m) -- (out);
  \draw[flow, draw=udemblue!70] (out.south) |- (fb.east);
  \draw[flow, draw=udemblue!70] (fb.north) -- (m.south);
\end{tikzpicture}

  \caption{Overview of \method{}. The method node is filled in Mila purple and the feedback path is drawn in UdeM blue, the two brand colours the class defines as \texttt{milapurple} and \texttt{udemblue}.}
  \label{fig:overview}
\end{figure*}

\subsection{Problem Setup}
\label{sec:problem-setup}

A displayed equation is part of the sentence around it and ends with that sentence's punctuation, as the attention operator of \citet{vaswani2017attention} does here,
\begin{equation}
  \operatorname{Attention}(Q, K, V) = \operatorname{softmax}\!\left(\frac{QK^\top}{\sqrt{d_k}}\right) V,
  \label{eq:attention}
\end{equation}
where $Q$, $K$ and $V$ are the query, key and value matrices and $d_k$ is the key dimension.
Only equations the text refers to carry a label, as \cref{eq:attention} does.

\paragraph{Run-in headings.}
Below \verb|\subsection|, a paragraph heading such as this one takes the place of \verb|\subsubsection|, which the lab does not use.
